%% LyX 2.3.6.2 created this file.  For more info, see http://www.lyx.org/.
%% Do not edit unless you really know what you are doing.
\documentclass[english]{beamer}
\usepackage[T1]{fontenc}
\usepackage[latin9]{inputenc}
\setcounter{secnumdepth}{3}
\setcounter{tocdepth}{3}
\usepackage{babel}
\usepackage{textcomp}
\usepackage{amstext}
\usepackage[all]{xy}
\ifx\hypersetup\undefined
  \AtBeginDocument{%
    \hypersetup{unicode=true,pdfusetitle,
 bookmarks=true,bookmarksnumbered=false,bookmarksopen=false,
 breaklinks=false,pdfborder={0 0 1},backref=false,colorlinks=true}
  }
\else
  \hypersetup{unicode=true,pdfusetitle,
 bookmarks=true,bookmarksnumbered=false,bookmarksopen=false,
 breaklinks=false,pdfborder={0 0 1},backref=false,colorlinks=true}
\fi

\makeatletter

%%%%%%%%%%%%%%%%%%%%%%%%%%%%%% LyX specific LaTeX commands.
%% Because html converters don't know tabularnewline
\providecommand{\tabularnewline}{\\}

%%%%%%%%%%%%%%%%%%%%%%%%%%%%%% Textclass specific LaTeX commands.
% this default might be overridden by plain title style
\newcommand\makebeamertitle{\frame{\maketitle}}%
% (ERT) argument for the TOC
\AtBeginDocument{%
  \let\origtableofcontents=\tableofcontents
  \def\tableofcontents{\@ifnextchar[{\origtableofcontents}{\gobbletableofcontents}}
  \def\gobbletableofcontents#1{\origtableofcontents}
}
\newenvironment{lyxcode}
  {\par\begin{list}{}{
    \setlength{\rightmargin}{\leftmargin}
    \setlength{\listparindent}{0pt}% needed for AMS classes
    \raggedright
    \setlength{\itemsep}{0pt}
    \setlength{\parsep}{0pt}
    \normalfont\ttfamily}%
   \def\{{\char`\{}
   \def\}{\char`\}}
   \def\textasciitilde{\char`\~}
   \item[]}
  {\end{list}}

%%%%%%%%%%%%%%%%%%%%%%%%%%%%%% User specified LaTeX commands.
\usetheme[secheader]{Boadilla}
\usecolortheme{seahorse}
\title[Theory of typeclasses]{Typeclasses from the viewpoint of type theory and category theory}
\subtitle{For functional programmers, with Haskell examples}
\author{Sergei Winitzki}
\date{2026-07-26}
\institute[ABTB]{Academy by the Bay}
\setbeamertemplate{headline}{} % disable headline at top
\setbeamertemplate{navigation symbols}{} % disable navigation bar at bottom
\usepackage[all]{xy} % xypic
%\makeatletter
% Macros to assist LyX with XYpic when using scaling.
\newcommand{\xyScaleX}[1]{%
\makeatletter
\xydef@\xymatrixcolsep@{#1}
\makeatother
} % end of \xyScaleX
\makeatletter
\newcommand{\xyScaleY}[1]{%
\makeatletter
\xydef@\xymatrixrowsep@{#1}
\makeatother
} % end of \xyScaleY

% Double-stroked fonts to replace the non-working \mathbb{1}.
\usepackage{bbold}
\DeclareMathAlphabet{\bbnumcustom}{U}{BOONDOX-ds}{m}{n} % Use BOONDOX-ds or bbold.
\newcommand{\custombb}[1]{\bbnumcustom{#1}}
% The LyX document will define a macro \bbnum{#1} that calls \custombb{#1}.

\usepackage{relsize} % make math symbols larger or smaller
\usepackage{stmaryrd} % some extra symbols such as \fatsemi
% Note: using \forwardcompose inside a \text{} will cause a LaTeX error!
\newcommand{\forwardcompose}{\hspace{1.5pt}\ensuremath\mathsmaller{\fatsemi}\hspace{1.5pt}}


% Make underline green.
\definecolor{greenunder}{rgb}{0.1,0.6,0.2}
%\newcommand{\munderline}[1]{{\color{greenunder}\underline{{\color{black}#1}}\color{black}}}
\def\mathunderline#1#2{\color{#1}\underline{{\color{black}#2}}\color{black}}
% The LyX document will define a macro \gunderline{#1} that will use \mathunderline with the color `greenunder`.
%\def\gunderline#1{\mathunderline{greenunder}{#1}} % This is now defined by LyX itself with GUI support.

\makeatother

\begin{document}
\global\long\def\gunderline#1{\mathunderline{greenunder}{#1}}%
\global\long\def\bef{\forwardcompose}%
\global\long\def\bbnum#1{\custombb{#1}}%
\frame{\titlepage}

\begin{frame}{Typeclasses: an algebraic view}

This material will appear in Chapter 13 of \emph{The Science of Functional
Programming}: \href{https://github.com/winitzki/sofp}{https://github.com/winitzki/sofp}

\vspace{2\baselineskip}

Plan:
\begin{enumerate}
\item Examine the typeclass type signatures:
\begin{itemize}
\item \textbf{algebraic} typeclasses ($P\,t\rightarrow t$)
\begin{itemize}
\item examples: \texttt{Monoid}, \texttt{Functor}, \texttt{Monad}, \texttt{Applicative} 
\end{itemize}
\item \textbf{co-algebraic} typeclasses ($t\rightarrow P\,t$)
\begin{itemize}
\item examples: \texttt{Show}, \texttt{Traversable}
\end{itemize}
\item the \texttt{Eq} typeclass is neither algebraic nor co-algebraic
\end{itemize}
\item Discover the universal patterns of typeclass derivation
\item Define typeclasses as categories (with \textbf{typeclass-preserving}
morphisms)
\item Express typeclass laws via (co)monad algebra laws
\item Prove the validity of the universal derivation patterns
\end{enumerate}
\end{frame}

\begin{frame}{Summary of typeclass derivation patterns}

{\footnotesize{}\vspace{-0.2cm}}{\footnotesize\par}
\begin{center}
{\footnotesize{}\vspace{-0.10cm}\hspace{-0.1cm}}%
\begin{tabular}{|c|c|c|}
\hline 
 & \textbf{\small{}Algebraic} & \textbf{\small{}Co-algebraic}\tabularnewline
\hline 
\hline 
\texttt{\small{}class C t where} & \texttt{\small{}m ::~P t $\rightarrow$ t} & \texttt{\small{}m ::~t $\rightarrow$ P t}\tabularnewline
\hline 
 & \texttt{\small{}instance C ()} & \texttt{\small{}instance C Void}\tabularnewline
\hline 
\texttt{\small{}instance (C a, C b)}{\small{} $\Rightarrow$} & \texttt{\small{}C (a, b)} & \texttt{\small{}C (Either a b)}\tabularnewline
\hline 
\texttt{\small{}instance (C a) $\Rightarrow$} & \texttt{\small{}C (r $\rightarrow$ a)} & \texttt{\small{}C (r, a)}\tabularnewline
\hline 
\texttt{\small{}instance (C x) $\Rightarrow$ C (F x)} & \texttt{\small{}C (Fix F)} & \texttt{\small{}C (Fix F)}\tabularnewline
\hline 
(\emph{for any} \texttt{t}) & \texttt{\small{}instance (Free t)} & \texttt{\small{}instance (Cofree t)}\tabularnewline
\hline 
\end{tabular}
\par\end{center}
\begin{itemize}
\item {\footnotesize{}\vspace{-0.0cm}}The new instances can be found for
all typeclasses \texttt{C}
\item The laws hold automatically for all new instances!
\item Remarks:
\begin{itemize}
\item $P$ and $F$ must be covariant functors (have lawful ``\texttt{fmap}''
methods)
\item \texttt{Free} and \texttt{Cofree} must be defined separately for each
typeclass \texttt{C}
\item Certain typeclasses \texttt{C} support other, specific derivation
patterns
\end{itemize}
\end{itemize}
\end{frame}

\begin{frame}{Typeclass method signatures: examples I}

\begin{itemize}
\item \vspace{-0.15cm}Gather all typeclass methods into one type signature
\item See if that signature is algebraic or co-algebraic
\end{itemize}
{\footnotesize{}\vspace{0.25cm}}{\footnotesize\par}
\begin{itemize}
\item Example: the \texttt{Monoid} typeclass
\begin{lyxcode}
class~Monoid~t~where

~~~~mappend~::~t~$\rightarrow$~t~$\rightarrow$~t

~~~~mempty~::~t
\end{lyxcode}
\begin{itemize}
\item uncurry the type of \texttt{mappend} as \texttt{(t, t) $\rightarrow$
t}
\item replace the type of \texttt{mempty} by \texttt{() $\rightarrow$ t}
\item combine into one function: \texttt{Maybe (t, t) $\rightarrow$ t}
\item this has the form \texttt{P t $\rightarrow$ t} , where \texttt{P
x = Maybe (x, x)}
\item so, the \texttt{Monoid} typeclass is \textbf{algebraic}
\end{itemize}
\item Terminology borrowed from category theory: 
\begin{itemize}
\item a ``$P$-functor algebra'' is a type $t$ with a function of type
$P\,t\rightarrow t$
\item a ``$P$-functor coalgebra'' is a type $t$ with a function of type
$t\rightarrow P\,t$
\item the typeclass methods are gathered into a single \textbf{structure
map}
\end{itemize}
\end{itemize}
\end{frame}

\begin{frame}{Typeclass method signatures: examples II}

{\footnotesize{}\vspace{-0.15cm}}The \texttt{Functor} typeclass
\begin{itemize}
\item The typeclass applies to type \emph{constructors}:
\begin{lyxcode}
class~Functor~f~where

~~~~fmap~::~(a~$\rightarrow$~b)~$\rightarrow$~f~a~$\rightarrow$~f~b
\end{lyxcode}
\begin{itemize}
\item uncurry the type: \texttt{(a $\rightarrow$ b, f a) $\rightarrow$
f b}
\item this has the form \texttt{P f b $\rightarrow$ f b}, where we define
\texttt{P} by:\\
 \texttt{P f x = (a $\rightarrow$ x, f a)}
\begin{itemize}
\item note that the type parameter ``\texttt{a}'' is used existentially
\end{itemize}
\end{itemize}
\item The \texttt{Functor} typeclass is \textbf{algebraic}
\item The \texttt{Functor} typeclass is also \textbf{co-algebraic}:
\begin{itemize}
\item flip the curried arguments: \texttt{f a $\rightarrow$ (a $\rightarrow$
b) $\rightarrow$ f b}
\item this is of the form \texttt{f a $\rightarrow$ P f a} , where we define
\texttt{P} by:\\
 \texttt{P f x = (x $\rightarrow$ b) $\rightarrow$ f b} 
\begin{itemize}
\item the type parameter ``\texttt{b}'' is used existentially
\item the type parameter ``\texttt{x}'' is in a \emph{covariant} position
at right
\end{itemize}
\end{itemize}
\item The (co)algebras are at the level of \emph{type constructors}
\begin{itemize}
\item Instead of $P\,t\rightarrow t$ we have $\forall x.\,P\,f\,x\rightarrow f\,x$
\end{itemize}
\end{itemize}
\end{frame}

\begin{frame}{Typeclass method signatures: examples III}

\begin{itemize}
\item {\footnotesize{}\vspace{-0.15cm}}The \texttt{Show} typeclass:
\begin{lyxcode}
class~Show~t~where

~~~~show~::~t~$\rightarrow$~String
\end{lyxcode}
\begin{itemize}
\item The type signature has the form $t\rightarrow P\,t$, where $P$ is
defined by:
\end{itemize}
\begin{lyxcode}
P~x~=~String~~~~~-{}-~a~constant~functor
\end{lyxcode}
\begin{itemize}
\item The \texttt{Show} typeclass is \textbf{co-algebraic} 
\end{itemize}
\item {\footnotesize{}\vspace{0.25cm}}The \texttt{Traversable} typeclass:
\begin{lyxcode}
class~(Applicative~g)~$\Rightarrow$~Traversable~f~where

~~~~traverse~::~(a~$\rightarrow$~g~b)~$\rightarrow$~f~a~$\rightarrow$~g~(f~b)
\end{lyxcode}
\begin{itemize}
\item flip the curried arguments: \texttt{f a $\rightarrow$ (a $\rightarrow$
g b) $\rightarrow$ g (f b)}
\item this has the form \texttt{f a $\rightarrow$ P f }a, where we define
\texttt{P} by:\\
 \texttt{P f x = (x $\rightarrow$ g b) $\rightarrow$ g (f b)}
\begin{itemize}
\item the type parameter ``\texttt{b}'' is used existentially
\item the type parameter ``\texttt{x}'' is in a \emph{covariant} position
at right
\end{itemize}
\item The \texttt{Traversable} typeclass is \textbf{co-algebraic}
\end{itemize}
\end{itemize}
\end{frame}

\begin{frame}{Typeclass method signatures: examples IV}

\begin{itemize}
\item {\footnotesize{}\vspace{-0.15cm}}The \texttt{Pointed} typeclass:
\begin{lyxcode}
class~Pointed~f~where

~~~~pure~::~t~$\rightarrow$~f~t
\end{lyxcode}
\begin{itemize}
\item The type signature has the form $P\:f\:x\rightarrow f\,x$, where
$P$ is defined by:
\end{itemize}
\begin{lyxcode}
P~f~x~=~x~~~~~-{}-~the~identity~functor
\end{lyxcode}
\begin{itemize}
\item The \texttt{Pointed} typeclass is \textbf{algebraic} 
\end{itemize}
\item {\footnotesize{}\vspace{0.25cm}}The \texttt{Eq} typeclass:
\begin{lyxcode}
class~Eq~t~where

~~~~equal~::~t~$\rightarrow$~t~$\rightarrow$~Bool
\end{lyxcode}
\begin{itemize}
\item The type signature is \emph{not} of the form $P\,t\rightarrow t$
or $t\rightarrow P\,t$ with a covariant functor $P$
\item The \texttt{Eq} typeclass is neither algebraic nor co-algebraic
\item The theory developed here will mostly not apply to \texttt{Eq}
\begin{itemize}
\item Typeclass derivation for \texttt{Eq} will be \emph{ad hoc}
\end{itemize}
\end{itemize}
\end{itemize}
\end{frame}

\begin{frame}{Typeclasses, ADTs, and OO interfaces}

\vspace{-0.15cm}ADT (algebraic data type)
\begin{lyxcode}
data~G~where

~~c1~::~Int~$\rightarrow$~G

~~c2~::~Bool~$\rightarrow$~G~$\rightarrow$~G
\end{lyxcode}
\begin{itemize}
\item Gather all constructors into a single type signature \texttt{P G $\rightarrow$
G}, where \texttt{P} is defined by \texttt{P x = Either Int (Bool,
x)}
\item \texttt{G} is a type that supports a value \texttt{m ::~P G $\rightarrow$
G}
\end{itemize}
OO (object-oriented) interface
\begin{lyxcode}
trait~G:~~~//~Scala~3

~~~def~f1~:~Boolean

~~~def~f2(other:~Int)~:~G
\end{lyxcode}
\begin{itemize}
\item Gather all methods into a single type signature \texttt{G $\rightarrow$
P G}, where \texttt{P} is defined by \texttt{P x = (Bool, Int $\rightarrow$
x)}
\item \texttt{G} is a type that supports a value \texttt{m ::~G $\rightarrow$
P G}
\end{itemize}
Typeclass
\begin{itemize}
\item Type \texttt{t} belongs to the typeclass if there is a value m ::~Q
t
\end{itemize}
\end{frame}

\begin{frame}{Derivation patterns for algebraic typeclasses I}

\begin{itemize}
\item {\footnotesize{}\vspace{-0.15cm}}Given \texttt{p ::~P a $\rightarrow$
a} and \texttt{q ::~P b $\rightarrow$ b}, what other types $c$
can we construct such that we can derive \texttt{r ::~P c $\rightarrow$
c}?
\begin{itemize}
\item Here \texttt{P} is viewed as an arbitrary fixed functor
\item We ignore the typeclass laws for now
\end{itemize}
\item For\texttt{ c = ()} , set \texttt{r} to the constant function of type
\texttt{P () $\rightarrow$ ()}
\item For \texttt{c = d $\rightarrow$ a} (where \texttt{d} is any fixed
type), write:
\end{itemize}
\begin{lyxcode}
r~::~P~(d~$\rightarrow$~a)~$\rightarrow$~d~$\rightarrow$~a

r~k~x~=~p~\$~fmap~(\textbackslash y~$\rightarrow$~y~x)~k
\end{lyxcode}
\begin{itemize}
\item For \texttt{c = (a, b)}, write:
\end{itemize}
\begin{lyxcode}
r~::~P~(a,~b)~$\rightarrow$~(a,~b)

r~k~=~(p~x,~q~y)~where

~~~~x~=~fmap~fst~k

~~~~y~=~fmap~snd~k
\end{lyxcode}
\begin{itemize}
\item But we cannot implement \texttt{r} for \texttt{c = Either a b} in
general
\begin{itemize}
\item because we cannot map \texttt{P (Either a b) $\rightarrow$ Either
(P a) (P b)}
\begin{itemize}
\item although this is certainly possible for some \texttt{P}
\end{itemize}
\end{itemize}
\end{itemize}
\end{frame}

\begin{frame}{Derivation patterns for algebraic typeclasses II}

{\footnotesize{}\vspace{-0.15cm}}Recursive types (described via pattern
functors)
\begin{itemize}
\item For a covariant functor $F$, the type $c$ is a \textbf{fixpoint}
of $F$ if:
\end{itemize}
\begin{lyxcode}
fix~::~F~c~$\rightarrow$~c

unfix~::~c~$\rightarrow$~F~c

fix~.~unfix~=~id~,~unfix~.~fix~=~id
\end{lyxcode}
\begin{itemize}
\item Assume that $F$ ``preserves the typeclass membership'': whenever
a type $t$ belongs to the typeclass then so does $F\:t$. This is
expressed via the function:
\end{itemize}
\begin{lyxcode}
h~::~(P~t~$\rightarrow$~t)~$\rightarrow$~P~(F~t)~$\rightarrow$~F~t

h~=~...~~~~~-{}-~assume~that~this~function~is~known
\end{lyxcode}
\begin{itemize}
\item Then \texttt{c} belongs to the typeclass:
\end{itemize}
\begin{lyxcode}
r~::~P~c~$\rightarrow$~c

r~=~fix~.~h~r~.~fmap~unfix
\end{lyxcode}
\end{frame}

\begin{frame}{Derivation patterns for algebraic typeclasses III}

{\footnotesize{}\vspace{-0.15cm}}Free typeclass instances
\begin{itemize}
\item A \textbf{free instance constructor} for an algebraic typeclass is
a functor $M$ such that:
\begin{itemize}
\item For any type $t$, the type $M\,t$ belongs to the typeclass
\begin{itemize}
\item The structure map is \texttt{m ::~P (M t) $\rightarrow$ M t}
\end{itemize}
\item For any type $c$ that belongs to the typeclass there is a unique
typeclass-preserving function of type $M\,c\rightarrow c$
\begin{itemize}
\item The ``typeclass-preserving'' property will be explained shortly
\end{itemize}
\item There are some other technical conditions, omitted for brevity
\end{itemize}
\item Examples:
\begin{itemize}
\item Free \texttt{Monoid} is the \texttt{List} type constructor
\item Free \texttt{Monad} is the (higher-kinded) type constructor defined
by:\\
\texttt{data FreeMonad f a = Pure a | Join f (FreeMonad f a)}
\end{itemize}
\item When a typeclass has no laws, the free instance is \texttt{FreeMonad
P}
\item Otherwise, free instance constructors must be defined \emph{ad hoc}
\end{itemize}
\end{frame}

\begin{frame}{Derivation patterns for co-algebraic typeclasses I}

\begin{itemize}
\item {\footnotesize{}\vspace{-0.15cm}}Given \texttt{p ::~a $\rightarrow$
P a} and \texttt{q ::~b $\rightarrow$ P b}, what other types $c$
can we construct such that we can derive \texttt{r ::~c $\rightarrow$
P c}?
\begin{itemize}
\item Here \texttt{P} is an arbitrary fixed functor
\item We ignore the typeclass laws for now
\end{itemize}
\item For\texttt{ c = Void} , set \texttt{r} to the function \texttt{absurd}
of type \texttt{Void $\rightarrow$ P Void}
\item For \texttt{c = (d, a)} (where \texttt{d} is any fixed type), write:
\end{itemize}
\begin{lyxcode}
r~::~(d,~a)~$\rightarrow$~P~(d,~a)

r~(k,~x)~=~fmap~(\textbackslash y~$\rightarrow$~(k,~y))~(p~x)
\end{lyxcode}
\begin{itemize}
\item For \texttt{c = Either a b}, write:
\end{itemize}
\begin{lyxcode}
r~::~Either~a~b~$\rightarrow$~P~(Either~a~b)

r~=~\textbackslash case

~~~Left~x~$\rightarrow$~p~x

~~~Right~y~$\rightarrow$~q~y
\end{lyxcode}
\begin{itemize}
\item But we cannot implement \texttt{r} for \texttt{c = (a, b)} in general
\begin{itemize}
\item because we cannot always map \texttt{(P a, P b) $\rightarrow$ P (a,
b)}
\begin{itemize}
\item although this is certainly possible for some \texttt{P}
\end{itemize}
\end{itemize}
\end{itemize}
\end{frame}

\begin{frame}{Derivation patterns for co-algebraic typeclasses II}

{\footnotesize{}\vspace{-0.15cm}}Recursive types (described via pattern
functors)
\begin{itemize}
\item For a covariant functor $F$, the type $c$ is a \textbf{fixpoint}
of $F$ if:
\end{itemize}
\begin{lyxcode}
fix~::~F~c~$\rightarrow$~c

unfix~::~c~$\rightarrow$~F~c

fix~.~unfix~=~id~,~unfix~.~fix~=~id
\end{lyxcode}
\begin{itemize}
\item Assume that $F$ ``preserves the typeclass membership'': whenever
a type $t$ belongs to the typeclass then so does $F\:t$. This is
expressed via the function:
\end{itemize}
\begin{lyxcode}
h~::~(t~$\rightarrow$~P~t)~$\rightarrow$~F~t~$\rightarrow$~P~(F~t)

h~=~...~~~~~-{}-~assume~that~this~function~is~known
\end{lyxcode}
\begin{itemize}
\item Then \texttt{c} belongs to the typeclass:
\end{itemize}
\begin{lyxcode}
r~::~c~$\rightarrow$~P~c

r~=~fmap~fix~.~h~r~.~unfix
\end{lyxcode}
\end{frame}

\begin{frame}{Derivation patterns for co-algebraic typeclasses III}

{\footnotesize{}\vspace{-0.15cm}}Cofree typeclass instances
\begin{itemize}
\item A \textbf{cofree instance constructor} for a co-algebraic typeclass
is a functor $N$ such that:
\begin{itemize}
\item For any type $t$, the type $N\,t$ belongs to the typeclass
\begin{itemize}
\item The structure map is \texttt{m ::~N t $\rightarrow$ P (N t)}
\end{itemize}
\item For any type $c$ that belongs to the typeclass there is a unique
typeclass-preserving function of type $c\rightarrow N\,c$
\item There are some other technical conditions, omitted for brevity
\end{itemize}
\item Examples:
\begin{itemize}
\item Cofree \texttt{Show} is:\\
\texttt{type CofreeShow t = (t, String)}
\item Cofree\texttt{ Comonad} is defined by (must use lazy pair!):\\
\texttt{data CofreeComonad f x = (x, f (CofreeComonad f x))}
\end{itemize}
\item When a typeclass has no laws, \texttt{CofreeComonad P} works
\item Otherwise, cofree instance constructors must be defined \emph{ad hoc}
\end{itemize}
\end{frame}

\begin{frame}{Typeclass-preserving functions: Algebraic case. I}

\begin{itemize}
\item \vspace{-0.15cm}Take an algebraic typeclass whose structure map has
type $P\,t\rightarrow t$. (Here $P$ is a covariant functor.)
\end{itemize}
\begin{lyxcode}
class~C~t~where

~~~~m~::~P~t~$\rightarrow$~t
\end{lyxcode}
\begin{itemize}
\item Consider two types $t$ and $u$ that belong to the typeclass
\item We say that a function $f::t\rightarrow u$ is \textbf{typeclass-preserving}
if $f$ always maps the operations on $t$ into the corresponding
operations on $u$
\item Example: the \texttt{Monoid} typeclass. A function $f$ is ``monoid-preserving''
if it maps both operations from the monoid $t$ to the monoid $u$
\item The operation \texttt{mappend}:
\begin{itemize}
\item The expression \texttt{mappend $t_{1}$ $t_{2}$} produces a new value
of type $t$
\item Apply $f$ to that value, obtain a value of type $u$
\item That value should be always equal to \texttt{mappend ($f$ $t_{1}$)
($f$ $t_{2}$)}
\end{itemize}
\item The operation \texttt{mempty}:
\begin{itemize}
\item The expression ``\texttt{mempty}'' is a value of type $t$
\item Apply $f$ to that value, obtain a value of type $u$
\item That value should be equal to \texttt{mempty} for the type $u$
\end{itemize}
\end{itemize}
\end{frame}

\begin{frame}{Typeclass-preserving functions: Algebraic case. II}

\vspace{-0.15cm}Formulate the ``typeclass-preserving'' property
in terms of $P$ alone
\begin{itemize}
\item Example: the \texttt{Monoid} typeclass with \texttt{P t = Maybe (t,
t)} and the structure map \texttt{$m_{t}$ ::~P t $\rightarrow$
t.}
\begin{itemize}
\item The function ``\texttt{mappend}'' corresponds to the structure map
$m_{t}$ applied to a value of the form \texttt{Just ($t_{1}$, $t_{2}$)}
\item The function ``\texttt{mempty}'' corresponds to $m_{t}$ applied
to \texttt{Nothing}
\item The substituted arguments in Just \texttt{($f$ $t_{1}$, $f$ $t_{2}$)}
correspond to applying (\texttt{fmap f}) to the value \texttt{Just
($t_{1}$, $t_{2}$)}
\item Applying (\texttt{fmap f}) to \texttt{Nothing} produces again \texttt{Nothing}
\end{itemize}
\item The typeclass-preserving condition becomes: for any \texttt{s} of
type \texttt{P t}: 
\[
m_{u}\,(\text{fmap}\,f\,s)=f\,(m_{t}\,s)\quad\text{or:}\quad f\circ m_{t}=m_{u}\circ\text{fmap }f
\]
\[
\xymatrix{\xyScaleY{2.8pc}\xyScaleX{5.5pc}P\,t\ar[r]\sp(0.5){\text{fmap }f}\ar[d]\sp(0.5){m_{t}} & P\,u\ar[d]\sp(0.5){m_{u}}\\
t\ar[r]\sp(0.5){f} & u
}
\]

\begin{itemize}
\item In category theory, this is called the $P$-\textbf{algebra morphism}
law
\end{itemize}
\end{itemize}
\end{frame}

\begin{frame}{Typeclass-preserving functions: Co-algebraic case}

\begin{itemize}
\item \vspace{-0.15cm}Take a co-algebraic typeclass with methods described
as $t\rightarrow P\,t$
\end{itemize}
\begin{lyxcode}
class~C~t~where

~~~~m~::~t~$\rightarrow$~P~t
\end{lyxcode}
\begin{itemize}
\item Let $t$, $u$ be two types belonging to this typeclass, with structure
maps $m_{t}$ and $m_{u}$
\item A function $f:t\rightarrow u$ is \textbf{typeclass-preserving} if:
\[
\text{fmap }f\circ m_{t}=m_{u}\circ f
\]
\[
\xymatrix{\xyScaleY{2.8pc}\xyScaleX{5.5pc}t\ar[r]\sp(0.5){f}\ar[d]\sp(0.5){m_{t}} & u\ar[d]\sp(0.5){m_{u}}\\
P\,t\ar[r]\sp(0.5){\text{fmap }f} & P\,u
}
\]
\item Example: a \texttt{Show}-preserving function; \texttt{P t = String}
\[
\text{show }(f\,x)=\text{show }x
\]
\item In other words, $f$ preserves the result of applying \texttt{show}
\end{itemize}
\end{frame}

\begin{frame}{Typeclass-preserving functions: General case I}

\begin{itemize}
\item \vspace{-0.15cm}Gather all methods of a typeclass into a profunctor
$Q$:
\end{itemize}
\begin{lyxcode}
class~C~t~where

~~~~m~::~Q~t~t
\end{lyxcode}
\begin{itemize}
\item Q is \emph{contravariant} in the first type parameter and \emph{covariant}
in the second; has \texttt{dimap ::~(a $\rightarrow$ b) $\rightarrow$
(c $\rightarrow$ d) $\rightarrow$ Q b c $\rightarrow$ Q a d}
\item Instead of a structure map, we have a typeclass instance value \texttt{m}
\item If types $t$ and $u$ belong to this typeclass, a function $f::t\rightarrow u$
is \textbf{typeclass-preserving} if the instance values $m_{t}$ and
$m_{u}$ satisfy the ``wedge condition'':
\[
\xymatrix{\xyScaleY{2.8pc}\xyScaleX{5.5pc}Q\,t\,t\ar[r]\sp(0.5){\text{dimap id }f} & Q\,t\,u\\
 & Q\,u\,u\ar[u]\sp(0.5){\text{dimap }f\text{ id}}
}
\]
\[
\text{dimap id }f\,m_{t}=\text{dimap }f\text{ id }m_{u}
\]
\end{itemize}
\end{frame}

\begin{frame}{Typeclass-preserving functions: General case II}

\begin{itemize}
\item \vspace{-0.15cm}As an example, take the Eq typeclass
\item Gather all methods of a typeclass into a \textbf{profunctor} $Q$:
\end{itemize}
\begin{lyxcode}
type~Q~a~b~=~(a,~a)~$\rightarrow$~Bool

dimap~::~(a~$\rightarrow$~b)~$\rightarrow$~(c~$\rightarrow$~d)~$\rightarrow$~Q~b~c~$\rightarrow$~Q~a~d

dimap~f~g~q~=~\textbackslash (x,~y)~$\rightarrow$~q~(f~x,~f~y)
\end{lyxcode}
\begin{itemize}
\item The wedge condition for $f::t\rightarrow u$ is: $m_{u}\,(f\,x,f\,y)=m_{t}\,(x,y)$
\item In other words, $f$ preserves the result of applying \texttt{equal}
\end{itemize}
\end{frame}

\begin{frame}{Typeclass-preserving functions: Category laws}

\vspace{-0.15cm}Typeclass-preserving functions are morphisms in a
category
\begin{itemize}
\item Identity functions ($\text{id}::t\rightarrow t$) are typeclass-preserving
\[
\text{dimap id id}=\text{id}
\]
\item If $f::t\rightarrow u$ and $g::u\rightarrow v$ are typeclass-preserving
then so is $g\circ f$
\[
\xymatrix{Q\,t\,t\ar[r]\sp(0.5){\text{dimap id }f}\ar[dd]\sp(0.5){\text{dimap id }(g\circ f)} & Q\,t\,u\ar[r]\sp(0.5){\text{dimap id }g} & Q\,t\,v\\
\xyScaleY{2.1pc}\xyScaleX{5.5pc} & Q\,u\,u\ar[u]\sp(0.5){\text{dimap }f\text{ id}}\ar[r]\sp(0.5){\text{dimap id }g} & Q\,u\,v\ar[u]\sp(0.5){\text{dimap }f\text{ id}}\\
Q\,t\,v &  & Q\,v\,v\ar[u]\sp(0.5){\text{dimap }g\text{ id}}\ar[ll]\sp(0.5){\text{dimap }(g\circ f)\text{ id}}
}
\]
\[
\text{dimap id }(g\circ f)\,m_{t}=\text{dimap }(g\circ f)\text{ id}\,m_{v}
\]
\end{itemize}
\end{frame}

\begin{frame}{Typeclass-preserving functions: Proof of composition law}

\vspace{-0.15cm}We have values $m_{t}$, $m_{u}$, $m_{v}$ such
that:
\[
\text{dimap id }f\,m_{t}=\text{dimap }f\text{ id }m_{u}\quad\text{and}\quad\text{dimap id }g\,m_{u}=\text{dimap }g\text{ id }m_{v}
\]
We need to show:
\[
\text{dimap id }(g\circ f)\,m_{t}\overset{?}{=}\text{dimap }(g\circ f)\text{ id}\,m_{v}
\]
Left-hand side:
\begin{align*}
\text{dimap id }(g\circ f)\,m_{t} & =\text{dimap id }g\,(\gunderline{\text{dimap id }f\,m_{t}})\\
 & =\text{dimap id }g\,(\text{dimap }f\text{ id }m_{u})=\text{dimap }f\,g\,m_{u}
\end{align*}
Right-hand side:
\begin{align*}
\text{dimap }(g\circ f)\text{ id}\,m_{v} & =\text{dimap }f\,\text{id }(\gunderline{\text{dimap }g\text{ id }m_{v}})\\
 & =\text{dimap }f\,\text{id }(\text{dimap id }g\,m_{u})=\text{dimap }f\,g\,m_{u}
\end{align*}

\end{frame}

\begin{frame}{Typeclass laws I: Example}

\vspace{-0.15cm}Typeclass laws are equations that the methods must
satisfy
\begin{itemize}
\item Example: monoid laws
\begin{align*}
 & \text{mappend}\,\text{mempty}\,x=x\quad,\quad\text{mappend}\,x\,\text{mempty}=x\quad,\\
 & \text{mappend}\,(\text{mappend}\,x\,y)\,z=\text{mappend}\,x\,(\text{mappend}\,y\,z)
\end{align*}
or in another notation:
\begin{align*}
 & e\oplus x=x\quad,\quad x\oplus e=x\quad,\\
 & (x\oplus y)\oplus z=x\oplus(y\oplus z)
\end{align*}
\item Generally, typeclass laws are one or more equations for the structure
map, using a number of arbitrary input values
\item It turns out that these equations can be replaced by more general
laws
\end{itemize}
\end{frame}

\begin{frame}{Typeclass laws II: monad algebras}

\vspace{-0.15cm}If $M$ is a monad with methods:
\begin{align*}
 & \text{pure}::a\rightarrow M\,a\\
 & \text{join}::M\,(M\,a)\rightarrow M\,a
\end{align*}
then an $M$-\textbf{monad algebra} is a type $t$ and a structure
map $m_{t}::M\,t\rightarrow t$ satisfying the following identity
and composition laws:
\[
m_{t}\circ\text{pure}=\text{id}\quad\quad\quad m_{t}\circ\text{join}=m_{t}\circ(\text{fmap }m_{t})
\]
\[
\xymatrix{t\ar[r]\sp(0.5){\text{pure}}\ar[rd]\sp(0.5){\text{id}} & M\,t\ar[d]\sp(0.45){m_{t}} & M\,(M\,t)\ar[r]\sp(0.5){\text{join}}\ar[d]\sp(0.45){\text{fmap }m_{t}} & M\,t\ar[d]\sp(0.45){m_{t}}\\
\xyScaleY{2.1pc}\xyScaleX{3.5pc} & t & M\,t\ar[r]\sp(0.5){m_{t}} & t
}
\]

\end{frame}

\begin{frame}{Typeclass laws III: deriving from monad algebras}

\vspace{-0.15cm}Often, typeclass laws can be expressed via monad
algebras

\vspace{0.15cm}Example: how to derive \texttt{Monoid}'s laws from
\texttt{List}-monad algebra laws
\begin{itemize}
\item The \texttt{List} monad: \texttt{pure x = {[}x{]}} and \texttt{join}
is list concatenation
\item Suppose \texttt{t} is a \texttt{List}-monad algebra, with \texttt{m
$::$ List t $\rightarrow$ t}
\begin{itemize}
\item Denote \texttt{e = m {[}{]}} and \texttt{x $\oplus$ y = m {[}x, y{]}}
\item The identity law says: \texttt{m {[}x{]} = x}
\item Use the composition law with argument \texttt{{[}{[}x{]}, {[}{]}{]}}
\begin{itemize}
\item \texttt{m (join {[}{[}x{]}, {[}{]}{]}) = m {[}x{]} = x}
\item \texttt{m (fmap m {[}{[}x{]}, {[}{]}{]}) = m {[}x, e{]} = x $\oplus$
e}
\end{itemize}
\item It follows that \texttt{x = x $\oplus$ e} and similarly \texttt{x
= e $\oplus$ x}
\item Use the composition law with arguments \texttt{{[}{[}x, y{]}, {[}z{]}{]}}
\begin{itemize}
\item \texttt{m (join {[}{[}x, y{]}, {[}z{]}{]}) = m {[}x, y, z{]}}
\item \texttt{m (fmap m {[}{[}x, y{]}, {[}z{]}{]}) = m {[}x $\oplus$ y,
z{]} = (x $\oplus$ y) $\oplus$ z}
\item So, \texttt{(x $\oplus$ y) $\oplus$ z = m {[}x, y, z{]} = x $\oplus$
(y $\oplus$ z)}
\end{itemize}
\item All other instances of the composition law will now hold
\end{itemize}
\item We have derived \texttt{Monoid}'s laws
\begin{itemize}
\item But we needed ingenuity to transform the laws into the familiar form
\end{itemize}
\end{itemize}
\end{frame}

\begin{frame}{Typeclass laws IV: algebraic typeclasses}

\begin{itemize}
\item \vspace{-0.15cm}If \texttt{t} is a \texttt{List}-monad algebra then
\texttt{t} is a \texttt{Monoid}
\item Generalize: for some monad \texttt{M}, if \texttt{t} is an \texttt{M}-monad
algebra then \texttt{t} belongs to some \emph{algebraic} typeclass
(with methods $m::P\,t\rightarrow t$)
\begin{itemize}
\item To which typeclass? (What functor is $P$?)
\item Given an algebraic typeclass, how to find a suitable monad \texttt{M}?
\end{itemize}
\item \textbf{Structure theorems} for algebraic typeclasses: 
\begin{enumerate}
\item If \texttt{M} is the \emph{free instance constructor} of a given algebraic
typeclass then \texttt{M} is a monad.
\item With that \texttt{M}: If \texttt{t} is an \texttt{M}-monad algebra
then \texttt{t} belongs to the typeclass.
\item If an algebraic typeclass has no laws then \texttt{M = FreeMonad P}
\end{enumerate}
\item It is not obvious how to find \texttt{M} given an algebraic typeclass
with a given structure functor \texttt{P} and a given set of laws
\begin{itemize}
\item The free instance constructor has to be derived separately for every
algebraic typeclass
\begin{itemize}
\item It is not clear that we will always find \texttt{M} within our type
system
\end{itemize}
\end{itemize}
\item It is also not obvious how to find \texttt{P} given \texttt{M}
\end{itemize}
\end{frame}

\begin{frame}{Typeclass laws V: why $M$ is always a monad}

\vspace{-0.15cm}Consider an algebraic typeclass with structure functor
$P$ and the free instance constructor $M$
\begin{itemize}
\item For a typeclass member $t$ we must have the ``evaluation map''
$M\,t\rightarrow t$
\begin{itemize}
\item This is a unique typeclass-preserving function of this type
\end{itemize}
\item Set $t=M\,u$ for any $u$, then we have a function $M\,(M\,u)\rightarrow M\,u$
\begin{itemize}
\item This is $M$'s \texttt{join} method
\end{itemize}
\item For any $t$ we must have a function $t\rightarrow M\,t$
\begin{itemize}
\item This is \texttt{$M$}'s \texttt{pure} method
\end{itemize}
\item The monad laws of $M$ follow from the properties of the evaluation
map (with some technical assumptions omitted here)
\end{itemize}
\end{frame}

\begin{frame}{Typeclass laws VI: laws and typeclass-preserving maps}

\vspace{-0.15cm}Consider an algebraic typeclass with structure functor
$P$ and the free instance constructor $M$
\begin{itemize}
\item Typeclass-preserving maps are $P$-algebra morphisms
\item All $M$-algebra morphisms are also $P$-algebra morphisms that obey
all laws, and vice versa
\end{itemize}
An algebraic typeclass without laws = the category of $P$-functor
algebras

An algebraic typeclass with laws = the category of $M$-monad algebras
\begin{itemize}
\item If $f::t\rightarrow u$ is a surjective $P$-functor algebra morphism,
and if a typeclass law holds in $t$, then that law also holds in
$u$
\item The evaluation map $M\,t\rightarrow t$ is a surjective typeclass-preserving
map
\begin{itemize}
\item It enforces the laws in $t$ as long as they hold in $M\:t$
\end{itemize}
\item If $f::t\rightarrow u$ is a typeclass-preserving map and some law
is violated in $t$ but holds in $u$ then the violations become invisible
after applying $f$
\end{itemize}
\end{frame}

\begin{frame}{Typeclass laws VII: co-algebraic typeclasses}

\vspace{-0.15cm}
\begin{itemize}
\item Co-algebraic typeclasses have methods $m::t\rightarrow P\,t$
\item The laws can be expressed via comonad coalgebras
\item \textbf{Structure theorems} for co-algebraic typeclasses: 
\begin{enumerate}
\item If \texttt{M} is the \emph{cofree instance constructor} of a given
co-algebraic typeclass then \texttt{M} is a comonad.
\item With that \texttt{M}: If \texttt{t} is an \texttt{M}-comonad coalgebra
then \texttt{t} belongs to the typeclass.
\item If a co-algebraic typeclass has no laws then \texttt{M = CofreeComonad
P}
\end{enumerate}
\item It is not obvious how to find \texttt{M} given a co-algebraic typeclass
with a given structure functor \texttt{P} and a given set of laws
\begin{itemize}
\item The cofree instance constructor has to be derived separately for every
co-algebraic typeclass
\item What is the cofree comonoid constructor? (I don't know.)
\end{itemize}
\item It is also not obvious how to find \texttt{P} given \texttt{M }
\end{itemize}
\end{frame}

\begin{frame}{Proving the algebraic typeclass derivation patterns I}

\vspace{-0.15cm}

The choice of the monad \texttt{M} encapsulates the full complexity
of the laws

Assume that \texttt{M} is given, then it remains to prove that:
\begin{enumerate}
\item \texttt{t = ()} is an \texttt{M}-monad algebra
\item If \texttt{t} is an \texttt{M}-monad algebra then so is \texttt{r
$\rightarrow$ t} for any fixed type \texttt{r}
\item If \texttt{a} and \texttt{b} are \texttt{M}-monad algebras then so
is \texttt{(a, b)}
\item If \texttt{F} is such that \texttt{F t} is an \texttt{M}-monad algebra
whenever \texttt{t} is, then any fixpoint of \texttt{F} is also an
\texttt{M}-monad algebra
\end{enumerate}
\end{frame}

\begin{frame}{Proving the algebraic typeclass derivation patterns II}

\vspace{-0.15cm}

Example: how to prove this for \texttt{u = r $\rightarrow$ t}
\begin{itemize}
\item Assume $m_{t}::M\,t\rightarrow t$ is a given $M$-monad algebra
\item We need to define $m_{u}::M\,(r\rightarrow t)\rightarrow r\rightarrow t$
\end{itemize}
\begin{lyxcode}
m\_u~::~M~(r~\textrightarrow{}~t)~$\rightarrow$~r~\textrightarrow{}~t

m\_u~mrt~x~=~m\_t~\$~fmap~(\textbackslash y~$\rightarrow$~y~x)~mrt
\end{lyxcode}
\begin{itemize}
\item Prove the $M$-monad algebra identity law for $m_{u}$, assuming that
the law already holds for $m_{t}$:
\end{itemize}
\begin{lyxcode}
(m\_u~.~pure)~f~x~=~m\_u~(pure~f)~x

~~~~~~~~~~~~~~~~~=~m\_t~\$~fmap~(\textbackslash y~$\rightarrow$~y~x)~(pure~f)

~~~~~~~~~~~~~~~~~=~m\_t~\$~pure~(f~x)

~~~~~~~~~~~~~~~~~=~f~x
\end{lyxcode}
\begin{itemize}
\item Prove the $M$-monad algebra composition law for $m_{u}$:
\begin{itemize}
\item (a longer derivation using naturality laws of $M$'s \texttt{join},
the functor composition laws, and the composition law of $m_{t}$)
\end{itemize}
\end{itemize}
\begin{lyxcode}
\end{lyxcode}
\end{frame}

\begin{frame}{Proving the co-algebraic typeclass derivation patterns}

\vspace{-0.15cm}

Assume that a comonad \texttt{M} is given, it remains to prove that:
\begin{enumerate}
\item \texttt{t = Void} is an \texttt{M}-comonad coalgebra
\item If \texttt{t} is an \texttt{M}-comonad coalgebra then so is \texttt{(r,
t)} for any fixed type \texttt{r}
\item If \texttt{a} and \texttt{b} are \texttt{M}-comonad coalgebras then
so is \texttt{Either a b}
\item If \texttt{F} is such that \texttt{F t} is an \texttt{M}-comonad coalgebra
whenever \texttt{t} is, then any fixpoint of \texttt{F} is also an
\texttt{M}-comonad coalgebra
\end{enumerate}
\end{frame}

\begin{frame}{Summary}
\begin{itemize}
\item Most typeclasses are algebraic or co-algebraic
\item Standard typeclass derivation combinators can be implemented as a
library
\begin{itemize}
\item The typeclass laws will hold automatically
\end{itemize}
\item Free (co-free) typeclass instances need to be implemented separately
\item Typeclass laws can be formulated via monad algebras or comonad coalgebras 
\end{itemize}
\end{frame}

\end{document}
